\begin{lemma}[Candidate-normalization sibling route]
\label{lem:metacic-confluence-audit-candidate-normalization-sibling-route}
Every accepted $\MetaCICConfluenceAuditPacketUp$ carrier that is read by a candidate-normalization consumer may cite \autoref{ch:concrete-instances-metacic-candidate-normalization-confluence-handoff-namecert} only as a boundary route. The audit packet contributes its conditional confluence rows, residual-substitution boundary, closed-face hypotheses, transport, replay, provenance, and local NameCert; the candidate-normalization handoff contributes the candidate gate, normal endpoint row, residual envelope, decidable boundary, and blocked-edge retention.
\end{lemma}
\begin{proof}
Project the audit packet through \autoref{def:metacic-confluence-audit-packet-carrier}. Its confluence-facing data are still the rows $D,C,N,A_0,S,R,H,T,P,L$, with $S$ retaining the substitution boundary and $R$ acting only as a separation pointer.

Read the sibling handoff as the neighboring candidate-normalization route of \autoref{ch:concrete-instances-metacic-candidate-normalization-confluence-handoff-namecert}. The composition identifies where the candidate-normalization consumer may attach to the audit map, but it does not add a new coordinate to the audit carrier.

Therefore the route preserves the existing no-claim boundary: it is not a proof of full Church--Rosser, subject reduction, unrestricted closed consistency, full substitution compatibility, or a premise-free parallel diamond.
\end{proof}
